\chapter{Conclusion}
\label{ch:conclusion}

The main goal of this thesis was to solve the tool fragmentation experienced by higher-education students and instructors. Rather than treating document viewing, communication, and AI assistance as isolated processes, this project aimed to develop a unified web application, StudySpace, that bridges the gap between passive content consumption and active collaboration. The resulting system serves as a functional proof of concept demonstrating the viability of this unified approach.

The resulting application successfully fulfils all functional requirements established at the start of the project. It provides Instructors with a course administration module to manage content and enrollments. For students, the system introduces a content extension system, allowing them to clone courses into a Private Workspace and submit Merge Proposals. The application also features a contextual messaging system that anchors discussions directly to source documents, and an integrated AI assistant that provides answers grounded exclusively in the uploaded Course Materials. 

By delivering these core capabilities within a secure, role-aware application, this thesis demonstrates that a unified workspace can effectively eliminate the need for disjointed learning and communication platforms.
% ----------------------------------------------------------------
\section{Future Work}
\label{sec:future-work}

While the current deployment successfully establishes the core workflows, it leaves several directions open for subsequent development.

The Contextual Anchor currently stores only the material identifier and title. A natural
extension would be to encode a page number or byte offset alongside the material reference,
allowing the receiving client to open the PDF viewer at the precise location the sender was
viewing at the time of writing.

The enrollment model supports two statuses (\texttt{ACTIVE} and \texttt{DROPPED}) but does not
include an enrollment request or Instructor approval step. A real faculty deployment would
require an invitation or explicit approval workflow before granting student access to Course
Materials in restricted courses.

Although the system already records cumulative study time and daily streak data per user,
this information is currently surfaced only in individual profile and dashboard views.
A practical extension would be to aggregate these metrics into a cohort-level engagement
overview, giving students a sense of their own progress relative to peers and providing
instructors with a lightweight tool for identifying students who may require additional
support. Such a feature would make fuller use of the data the system already collects
and could serve as a concrete motivational signal without introducing any change to the
existing database schema.

Finally, the backend was deployed to a single-instance PaaS environment. Achieving true
horizontal scalability under concurrent load would require two infrastructure upgrades:
externalising the WebSocket session state through a shared STOMP broker (such as RabbitMQ or
ActiveMQ), and provisioning the JWT signing secret through a dedicated secrets vault rather
than a plain environment variable.
